% Generated draft: requires mathematical and source review.
\documentclass[11pt]{article}
\usepackage[margin=25mm]{geometry}
\usepackage{fontspec}
\setmainfont{Latin Modern Roman}
\usepackage{amsmath,amssymb,mathtools,unicode-math}
\setmathfont{Latin Modern Math}
\usepackage{enumitem,needspace,xurl,hyperref}
\setlength{\parindent}{0pt}
\setlength{\parskip}{5pt}
\setlength{\emergencystretch}{3em}
\newcommand{\N}{\mathbb{N}}
\newcommand{\Z}{\mathbb{Z}}
\newcommand{\Q}{\mathbb{Q}}
\newcommand{\R}{\mathbb{R}}
\newcommand{\C}{\mathbb{C}}
\newcommand{\F}{\mathbb{F}}
\newcommand{\A}{\mathbb{A}}
\newcommand{\PP}{\mathbb{P}}
\newcommand{\E}{\mathbb{E}}
\newcommand{\eps}{\varepsilon}
\newcommand{\dd}{\,\mathrm{d}}
\begin{document}
\section*{UnsolvedMath ID: 10400170}
\textbf{Source problem code:} catalogue entry
\subsection{Definitions and mathematical statement}
Let \(\mathcal C_1\) be a semisimple ribbon category with finitely many simple isomorphism classes. Its braiding and ribbon trace give an \(S\)-matrix, whose entries, up to normalization, are traces of the double braidings on pairs of simple objects. Assume this matrix is not invertible. A modular extension here means a modular category \(\mathcal C_2\) containing \(\mathcal C_1\) as a full subcategory with its relevant tensor, braiding and ribbon structures; fullness means that morphism sets between objects of \(\mathcal C_1\) agree with those in \(\mathcal C_2\).

Y. Kawahigashi asks whether there always exists a \emph{minimal} such extension and whether it has an appropriate \emph{essential uniqueness}, as claimed in Theorem 2.13 of reference [315]. If these assertions hold, determine the relation between the two TQFTs arising from \(\mathcal C_1\) and its minimal modular extension. The source reports a quantum-double construction of an extension, citing [315, 316, 182], while allowing that other extensions can exist.

The excerpt does not define minimality or essential uniqueness. In particular, it supplies neither a dimension-minimization condition nor a universal property, nor an equivalence notion respecting the embedding. It also leaves the TQFT construction for the possibly degenerate category \(\mathcal C_1\) unspecified. These choices and the exact categorical hypotheses are required before the existence, uniqueness and TQFT-comparison questions become fixed formal assertions.
\subsection{Short English statement}
Does every category in the stated class with degenerate \(S\)-matrix have an essentially unique minimal modular extension? If so, compare the TQFT of the original category with that of the extension.
\subsection{Sources}
\begin{itemize}
\item Y. Kawahigashi, attributed Problem 9.6 in T. Ohtsuki, \emph{Problems on invariants of knots and 3-manifolds} (2002), supplied PDF page 140. \url{https://msp.org/gtm/2002/04/gtm-2002-04-024s.pdf}. Recorded access 2026-07-29; extraction pdf-font-structure.
\item Source references [315, 316, 182] concern the quantum-double extension; Theorem 2.13 of [315] is explicitly named for the minimality and uniqueness claim. Complete bibliography entries are absent.
\item Local UnsolvedMath registry, \texttt{data/unsolvedmath/problems.json}, permanent ID \(10400170\), exact source code \texttt{catalogue entry}. Its supplied summary mentions a M\"uger theorem and modular-extension theory as partial answers, without an exact theorem statement or additional reference.
\end{itemize}
\end{document}
